\begin{claim}{F.10}{Proved}{Discrete procyclic directions in the residual centralizer}
Let $A_5\leq\Gamma^+_{0,12}$ be the geometric
icosahedral group acting on its twelve-vertex orbit,
and put $P=P\Gamma_{0,12}$.  For every nonidentity
$x\in P$, its procyclic closure is a full
$\widehat\Z$, and
\[
 \overline{\langle x\rangle}\cap
              C_{\widehat P}(A_5)=1.                  \tag{F.10.1}
\]
In particular, the full procyclic subgroup of the
common forgetting kernel constructed in F.9 meets
the actual residual $A_5$ centralizer trivially.
For that particular word $w$, it already has
trivial intersection with the centralizer of the
specified involution $\beta=(1\ 2)(4\ 5)$ in $A_5$.
\end{claim}
\begin{proof}
\textbf{An explicit separating family of integral representations.}
We first construct the finite quotients used in the
proof.  For $n\geq4$, normalizing the last marked
point to infinity and the first two to $0,1$
identifies
\[
 M_{0,n}=\operatorname{Conf}_{n-3}
                              (\mathbf C\setminus\{0,1\}).
\]
The exact homeomorphism
\[
\begin{split}
 \operatorname{Conf}_{n-1}(\mathbf C)
   &\longrightarrow\mathbf C\times\mathbf C^*\times M_{0,n},\\
 (z_0,z_1,z_2,\ldots,z_{n-2})
   &\longmapsto\left(z_0,z_1-z_0,
        \left(\frac{z_j-z_0}{z_1-z_0}\right)_{j=2}^{n-2}\right)
\end{split}                                                   \tag{F.10.2}
\]
has inverse $(a,b,(u_j))\mapsto(a,a+b,(a+bu_j))$.
The section fixing the first two factors at $(0,1)$
therefore gives an injective homomorphism
\[
          P\Gamma_{0,n}\longrightarrow P_{n-1},
 \qquad P_{n-1}\cong\Z\times P\Gamma_{0,n}.             \tag{F.10.3}
\]
Here $P_{n-1}$ is the ordered plane braid group.
All basepoints in (F.10.2) are the displayed normalized
configurations.  Thus (F.10.3) is the actual direct-factor
map, rather than an inference from the ranks of the groups.

Put $r=n-1$ and $F_r=\langle t_1,\ldots,t_r\rangle$.
Artin's faithful braid action gives
$P_r\hookrightarrow\operatorname{Aut}(F_r)$
\cite[Theorem~14]{A-Artin1947}.  We use the homomorphism
convention
\[
 \sigma_i(t_i)=t_it_{i+1}t_i^{-1},\qquad
 \sigma_i(t_{i+1})=t_i,
\]
fixing the other generators, with automorphism products
composed in that order.  Equivalently, for chronological
point-motion concatenation, use the inverse of the right
monodromy action on the punctured-plane group.  This
conversion makes $A(bc)=A(b)\circ A(c)$ and preserves
faithfulness.  Every pure braid fixes $F_r^{\rm ab}$.

Let
\[
 \mu:F_r\longrightarrow
          \Z\langle\!\langle X_1,\ldots,X_r\rangle\!\rangle^*,
 \qquad \mu(t_i)=1+X_i
\]
be the integral Magnus expansion, with
$(1+X_i)^{-1}=1-X_i+X_i^2-\cdots$.
Here is the injectivity needed below.  Write a nonidentity
reduced word in maximal syllables as
$t_{i_1}^{a_1}\cdots t_{i_s}^{a_s}$, where every
$a_j\ne0$ and successive $i_j$ differ.  The coefficient
of the monomial $X_{i_1}\cdots X_{i_s}$ is
$a_1\cdots a_s\ne0$.  Indeed that monomial has $s$
successive runs of distinct variables.  Omitting a
syllable contribution leaves at most $s-1$ runs,
whereas using all syllables and total degree $s$
forces their linear terms.  Thus $\mu$ is injective.

For each $d\geq1$ let
\[
 T_{r,d}=\Z\langle X_1,\ldots,X_r\rangle/
                   (\text{monomials of degree }>d).
\]
An automorphism $\alpha$ of $F_r$ acts on this finite
free $\Z$-module by the substitution
\[
             X_i\longmapsto\mu(\alpha(t_i))-1
                                      \pmod{\deg>d}.    \tag{F.10.4}
\]
The substitution identities for group words show
directly that these are homomorphisms of automorphism
groups.  For an automorphism fixing $F_r^{\rm ab}$,
the right side is $X_i$ plus terms of degree at least
two.  Order the monomial basis by decreasing degree.
The matrix in (F.10.4) is then upper unitriangular,
and all its entries are integers.  Combining
(F.10.3)--(F.10.4) gives actual homomorphisms
\[
 \rho_d:P\Gamma_{0,n}\longrightarrow
                      \operatorname{UT}_{N_d}(\Z),
 \qquad N_d=1+r+\cdots+r^d.                             \tag{F.10.5}
\]
This family separates elements.  Two different pure
mapping classes give different automorphisms of $F_r$
by (F.10.3) and Artin faithfulness.  Some generator has
different images; injectivity of $\mu$ makes their
expansions different in some finite degree.  That
degree gives a representation (F.10.5) which separates
the two classes.  In particular, no passage from a
discrete lower-central quotient to a completed one
is assumed here: the detecting maps are integral
matrices and their reductions in finite matrix groups.

\textbf{The actual finite-group action on the matrices.}
Return to $n=12$ and put $G=A_5$.  Its intersection
with $P$ is trivial, as is seen from its faithful
puncture permutation.  The subgroup $PG$ of
$\Gamma^+_{0,12}$ is the actual semidirect product
$P\rtimes G$ and has finite index.  Thus its profinite
completion is the corresponding open subgroup of
$\widehat\Gamma^+_{0,12}$, with the specified copies
of $\widehat P$ and $G$.

The discrete centralizer $C_P(G)$ is trivial.
Indeed the quotient sphere has exactly the three
branch/puncture decorations of orders $2,3,5$.
The discrete covering normalizer sequence of F.7
has trivial quotient mapping class group, so its
centralizer is $Z(A_5)=1$.
For $x\ne1$ choose $a\in G$ with
$axa^{-1}\ne x$, and then choose $d$ in (F.10.5)
which separates these two elements.  Write $\rho=\rho_d$.
On a direct sum of $|G|$ copies of $\Z^{N_d}$,
indexed by $g\in G$, define
\[
       R(h)=\operatorname{diag}_{g\in G}
                              \rho(g^{-1}hg),\qquad h\in P.
                                                               \tag{F.10.6}
\]
Let $T_a$ carry the block with index $g$ to the block
with index $ag$.  Then
\[
 T_aT_b=T_{ab},\qquad
 T_aR(h)T_a^{-1}=R(aha^{-1}).                            \tag{F.10.7}
\]
For example the block with index $g$ on the left of
the second identity is
$\rho((a^{-1}g)^{-1}h(a^{-1}g))$
$=\rho(g^{-1}aha^{-1}g)$, which is its right-hand
block.  Consequently
\[
       h a\longmapsto R(h)T_a                          \tag{F.10.8}
\]
is a homomorphism from the actual group $P\rtimes G$
to $\operatorname{GL}_{|G|N_d}(\Z)$.
Its reductions modulo every positive integer are
finite quotients of this group.

\textbf{Every prime coordinate is detected.}
Put $W=R(x)$.  It is unipotent.  By the choice of
$d$, $T_aWT_a^{-1}\ne W$, already on the block
with index $1$.  The finite polynomial logarithm
\[
        L=\log W=\sum_{j\geq1}
                      \frac{(-1)^{j+1}}j(W-I)^j
\]
is a rational nilpotent matrix; the sum terminates.
The exponential and logarithm polynomials are
mutually inverse on nilpotent and unipotent matrices
over a field of characteristic zero.  Hence
\[
                       T_aLT_a^{-1}-L\ne0.             \tag{F.10.9}
\]
In particular this is the same nonzero rational
matrix after embedding $\Q$ in $\Q_p$, for every
prime $p$.

The reductions of (F.10.8) modulo $p^k$, for all
$k$, give a continuous homomorphism from
$\widehat{P\rtimes G}$ to
$\operatorname{GL}_{|G|N_d}(\Z_p)$.
The image of the procyclic power $x^u$, for
$u\in\widehat\Z$, is
\[
                          \exp(u_p L),                 \tag{F.10.10}
\]
where $u_p$ is its $\Z_p$ coordinate.
To check this equality, first take integral $u$;
the identity $W^u=\exp(uL)$ follows from the finite
polynomial identities.  For arbitrary $u_p$, pass
to limits of integers in $\Z_p$.  The matrix
polynomial is continuous, and the powers of $W$
modulo $p^k$ lie in an upper unitriangular finite
$p$-group, so these powers depend only on $u_p$.
The limit lies in $\operatorname{GL}(\Z_p)$ because
the integral matrices form a closed subset there.
This proves (F.10.10) through the specified finite
quotients.

Now an element $x^u$ centralizing $G$ commutes with
$a$.  Apply the continuous representation and then
the polynomial logarithm to (F.10.10).  It gives
\[
                   u_p\bigl(T_aLT_a^{-1}-L\bigr)=0
                      \quad\hbox{in }M_{|G|N_d}(\Q_p).
                                                               \tag{F.10.11}
\]
By (F.10.9), $u_p=0$ for every prime $p$.
Thus $u=0$ in $\widehat\Z$ and (F.10.1) follows.
The same argument without $T_a$, using any
$\rho_d(x)\ne I$, proves that
$u\mapsto x^u$ is injective.  Its image is exactly
the procyclic closure, which is therefore a full
$\widehat\Z$ as claimed.

\textbf{The particular F.9 direction.}
The following transport fixes the basepoint convention.
Put $\varphi=(1+\sqrt5)/2$, $\rho=\sqrt{\varphi+2}$,
and $z_*=-i(\rho+\varphi)$.  In the tetrahedral coordinate
of F.9 the orbit of $z_*$ is the specified icosahedral
twelve-point configuration.  The extra involution has
M\"obius matrix
\[
 \begin{pmatrix}
  1&\varphi-1-i\varphi\\
  \varphi-1+i\varphi&-1
 \end{pmatrix}.                                         \tag{F.10.12}
\]
This matrix represents $(1\ 2)(4\ 5)$ under the
fixed five-letter identification of $A_5$, and its
twelve-vertex permutation is checked in the certificate
below.  Use the
polygonal source path
$\gamma:z_*\longrightarrow5+z_*\longrightarrow b=5+4i$.
It misses $R_A$: the finite points of $R_A$ have real
and imaginary coordinates of absolute value less than
$3/2$, the first segment has imaginary part less than
$-3$, and the second has real part $5$.
Transport along the full labelled orbit path
$(g\gamma)_{g\in A_4}$ identifies the marked mapping
class groups at these two configurations.  Equivariant
isotopy extension keeps the specified $A_4$ action
fixed.  The symbol $\beta$ at the F.9 basepoint denotes
the transport of (F.10.12) by this path.  The exact
matrix and vertex-label certificate
\path{code/part_a/second_a4_c3_attack.py}
checks all $180$ generator products and all $720$
vertex images in the number field
$\Q(i,\sqrt5,\rho)$.

The F.9 source homomorphism is the full twelve-point
orbit map of F.7.8, before the stated forgetting.
For every source path, its configurations remain
invariant under the same fixed M\"obius $A_4$.
Equivariant isotopy extension along the path, or
equivalently the point-pushing lift of the quotient
cover in F.7.8, places its monodromy in the actual
discrete $C_P(A_4)$.  That map identifies the source
$F_{13}$ with this discrete centralizer.  Thus the
word $w$ is a nonidentity element of $C_P(A_4)$,
not merely an element whose seven-point images
are compatible with $A_4$.

In the fixed five-letter description, $A_4$ fixes
the fifth letter and
$\beta=(1\ 2)(4\ 5)$ carries it to the subgroup
fixing the fourth.  These subgroups generate $A_5$;
equivalently $\langle A_4,\beta\rangle=A_5$.
Therefore $\beta w\beta^{-1}\ne w$, since the
opposite equality would put $w$ in the just proved
trivial discrete $A_5$ centralizer.  We may take
$a=\beta$ throughout (F.10.6)--(F.10.11).
This proves the last assertion of the claim, on
the whole full profinite procyclic closure.

\textbf{Executed formula certificate.}
The exact-integer program
\path{code/part_a/procyclic_magnus_f10.py}
checks the maximal-syllable coefficient on all $484$
nonempty reduced words of length at most five in two
generators.  It verifies the substitution-composition
identities and upper unitriangular form for two actual
pure Artin automorphisms on the degree-three tensor
algebra of rank $40$, and checks eleven integral
power/logarithm identities.  For the actual action of
$A_5$ on $\operatorname{Inn}(F_5)$ by permutation of
the five free generators, it checks all $180$
displayed block-conjugation identities and $60$
block-product identities.  Finally it checks finite
power reductions modulo $p^2$ for $p=2,3,5,7,11$.
This is a certificate for the algebraic formulas;
it does not compute a separating tensor degree for
the particular F.9 word.  Existence of that degree
is proved above by faithful Artin action and the
integral coefficient argument.

The conclusion restricts the actual residual
centralizer on every displayed discrete procyclic
direction.  It does not assert that an arbitrary
element of that closed residual group lies in one
of these subgroups.  No such covering assertion is
used, and the full equality in F.7 remains unresolved.
\end{proof}
